\section{若干补充说明}\label{sec:misc}

\subsection{规范场在流时间零点附近的行为}

在正规化的理论中，随时间变化的场 $B_{\mu}(t,x)$ 满足边界条件 \eqref{eq:2.3}，
因此当 $t\to0$ 时其关联函数趋于裸规范场 $A_{\mu}(x)$ 的关联函数。
然而，由于后者需要用一个发散常数进行重正化，
在重正化并移除正规化之后，关联函数在 $t=0$ 处趋于奇异。
其 $t\to0$ 的渐近行为由一个按带奇异系数的定域重正化场展开的式子
\begin{equation}
   B_{\mu}(t,x)=c_B(t)(\ren{A})_{\mu}(x)+\rmO(t)
   \label{eq:8.1}
\end{equation}
描述 \cite{Symanzik}。例如在微扰论单圈阶，系数 $c_B(t)$ 在 $t=0$
处对数发散。

不难证明重正化群方程
\begin{equation}
   \left\{\mu{\partial\over\partial\mu}+\beta{\partial\over\partial g}
   -2\gamma\lambda{\partial\over\partial\lambda}
   +\gamma+\beta/g
   \right\}c_B(t)=\rmO(t)
   \label{eq:8.2}
\end{equation}
成立，其中
\begin{equation}
   \beta=-b_0g^3+\rmO(g^5),
   \qquad
   \gamma=c_0g^2+\rmO(g^4)
   \label{eq:8.3}
\end{equation}
分别是 $\beta$ 函数与规范场的反常量纲。例如在朗道规范下，
该方程容易积分，得到
\begin{equation}
   B_{\mu}(t,x)
   \mathrel{\mathop\sim_{t\to0}}
   \bigl(2b_0\gbar(q)^2\bigr)^{1/2-c_0/2b_0}
   R(g)(\ren{A})_{\mu}(x),
   \label{eq:8.4}
\end{equation}
其中 $\gbar(q)$ 是动量 $q=(8t)^{-1/2}$ 处的跑动耦合，
$R(g)$ 是把重正化场与重正化群不变的规范场联系起来的因子。

式 \eqref{eq:8.4} 是正规化理论中规范场所满足的边界条件的残余。
它表明由流方程生成的场以一种普适的方式与基本场相联系，
也就是说这里没有有限重正化的余地，原因在于任何这样的重正化
都会破坏 BRS 对称性。

\subsection{规范不变的复合场}

Wilson 圈与规范不变的定域场不依赖流方程中的参数 $\alpha_0$，
并且在正流时间处不需要重正化。在小流时间处，
它们的渐近行为由 $\SUn$ 规范理论中相应重正化场的标度性质决定，
从而反映出后者奇异的本质。

作为示例，考虑密度
\begin{equation}
   E=\frac{1}{4}G_{\mu\nu}^aG_{\mu\nu}^a,
   \label{eq:8.5}
\end{equation}
文献~\cite{WilsonFlow} 此前已对其作过研究。既然 $E$
可以与单位场混合，小流时间处的主导项是
\begin{align}
   E(t,x)=\langle E(t,x)\rangle+c_E(t)
   \ren{\Bigl\{\frac{1}{4}F^a_{\mu\nu}F^a_{\mu\nu}\Bigr\}}(x)+\rmO(t),
   \label{eq:8.6}
\end{align}
即其期望值，而第一个次领头项正比于基本场的重正化作用量密度。
相关系数 $\langle E(t,x)\rangle$ 与 $c_E(t)$ 满足一个重正化群方程，
由它可以解析地算出二者在小 $t$ 处精确的渐近行为。

\subsection{含物质场的理论}

存在物质场时，梯度流由与纯规范理论相同的方程定义
（小节 \ref{ssec:defflow}）。
至于把物质场也纳入时间演化这一有趣却更为复杂的情形，
本文不予考虑。

在可重正化理论中，物质场是量纲为 $1$ 的标量场或量纲为 $3/2$
的费米子场。只要正规化保持规范对称性，
我们在第 \ref{sec:brs}、\ref{sec:fin} 节中的论证便可逐字照搬。
特别地，涉及物质场的附加边界抵消项被整体对称性以及
其量纲不得超过 $4$ 的要求排除。
因此正流时间处关联函数的有限性同样得到保证。

QCD 是具备保证有限性的全部所需性质的规范理论的突出例子；
许多其他有兴趣的理论也是如此，
包括 SU(2) Higgs 模型、超对称版本的 QCD 以及人工色（technicolour）理论。

\subsection{格点正规化}

尽管费曼规则比维数正规化情形复杂，
预计梯度流的格点微扰分析不会遇到根本性的困难。

格点流方程中梯度项的一个可能选择是 Wilson 作用量的梯度
\cite{WilsonFlow}。然而没有理由非选这个特殊的作用量不可，
也没有理由要求它与理论本身的规范作用量一致。
只要该项在经典连续极限下（即在微扰论的树级）
具有正确的形式与归一化，
随时间变化的规范场的关联函数就不会依赖于所作的具体选择，
除了按格距 $a$ 正幂次消失的格点效应之外。

连续洛伦兹对称性在格点上破缺，
但剩余的严格对称性足以排除那些未曾在连续理论中出现过的抵消项。
因此我们预期：随时间变化的场的关联函数不需要重正化，
而且它们在连续极限下的取值与正规化无关
（耦合与规范固定参数的有限重正化除外）。
此外，若格点理论是 $\rmO(a)$ 改进的，
则关联函数自动也是改进的，
因为不存在同时具备所有必需性质的 $\rmO(a)$ 抵消项候选者。
